%!TEX root = ../template.tex

%% ----------------------------------------------------------------------------
%%  NOVATHESIS — 3-BackMatter/app-migration.tex
%% ----------------------------------------------------------------------------
%%
%% Version 8.0.1 (2026-08-05)
%% Copyright (C) 2004-26 by João M. Lourenço <joao.lourenco@fct.unl.pt>
%%
%% Migration guide: upgrading a document written for NOVAthesis 7.10.x or
%% earlier to 8.0, whose glossaries are driven by bib2gls.
%%
%% ----------------------------------------------------------------------------

\typeout{NT FILE 3-BackMatter/app-migration.tex}%

\chapter{Migrating from \ntindex[migration]{7.10.x}}
\label{app:migration}

\novathesis~8.0 changed how glossaries, acronyms and symbols are stored.  If
you are starting a new document, you can skip this appendix entirely --- the
template already ships in the new format.  If you are upgrading a document
you have been writing against 7.10.x or earlier, this is the only change that
requires your attention.

%%=======================================================================
\section{What Changed, and Why}
%%=======================================================================

Glossary entries used to be written as \hologo{LaTeX} commands in
\texttt{.tex} files, and processed by \texttt{makeglossaries}.  They are now
\texttt{.bib} entries, processed by \ntindex[packages!bib2gls]{\texttt{bib2gls}}.

The reason is a hard limit in \hologo{pdfTeX} and \hologo{XeTeX}: a document
can only write to 16 files at once, and the old glossary machinery claimed
four of those slots.  The template had to load the \texttt{morewrites}
package to work around it, which made every \hologo{pdfLaTeX} pass roughly
nine times slower.  With \texttt{bib2gls} the glossaries need \emph{no} write
slots at all, so most documents no longer need \texttt{morewrites}:
a full \hologo{pdfLaTeX} build of this manual went from 109 to 36 seconds.
See Section~\ref{sec:build_fastwrites}.

\begin{warningbox}[title={New requirement: Java}]
  \texttt{bib2gls} is a Java program.  It ships with \TeX~Live, but it needs a
  \gls{JRE} on your machine.  Overleaf provides one.  For a local
  installation, install any current \gls{JRE} if \verb!bib2gls --version!
  does not work.
\end{warningbox}

%%=======================================================================
\section{The Migration, Step by Step}
\label{sec:migration_steps}
%%=======================================================================

\paragraph{1. Convert your entry files.}
\begin{bashcode}
make glsbib
\end{bashcode}

This converts every glossary \texttt{.tex} file in \verb!1-FrontMatter/! into
a \texttt{.bib} file of the same name, leaving the originals untouched.  It
skips any file already converted, so it is safe to run twice.

\paragraph{2. Re-add your \texttt{sort} keys.}  This is the step people miss,
so \verb!make glsbib! counts them for you and prints a warning like:

\begin{verbatim}
  convert  1-FrontMatter/symbols.tex  ->  1-FrontMatter/symbols.bib
     WARNING: 1-FrontMatter/symbols.tex had 8 sort key(s). convertgls2bib
              DROPS them; re-add them in 1-FrontMatter/symbols.bib or those
              entries will mis-sort.
\end{verbatim}

The converter silently discards every \texttt{sort} key.  Nothing breaks ---
the document still builds, at the same page count --- but the affected
entries come out in the wrong order.  It bites exactly the entries that
needed a sort key in the first place: symbols, and any entry whose name is a
command.  Open each converted file and restore them:

\begin{latexcode}
@symbol{mu,
  name        = {\ensuremath{\mu}},
  sort        = {mu},        % <- restore this
  description = {Mu},
  type        = {symbols}
}
\end{latexcode}

\paragraph{3. Convert symbols to \texttt{@symbol}.}  The converter turns
everything into \texttt{@entry}.  Entries that are symbols should use
\texttt{@symbol} instead, so \texttt{bib2gls} treats them as symbols rather
than words.

\paragraph{4. Remove \latexinline{\glsaddall}.}  If your document calls
\latexinline{\glsaddall}, delete the call.  It has no equivalent in
\texttt{bib2gls}, which decides inclusion per glossary: the template already
lists every entry of the glossary and symbols lists, and only the acronyms
you actually use.

\paragraph{5. Delete the old \texttt{.tex} entry files} once the output looks
right.  They are no longer read.  If one is still registered in
\verb!0-Config/4_files.tex!, the build stops with an error telling you so ---
that error is the safety net for this whole procedure, not a bug.

\paragraph{6. Rebuild and check the lists.}  Compare the glossary, acronym
and symbol lists against your previous \gls{PDF}: same entries, same order.

%%=======================================================================
\section{If Something Looks Wrong}
%%=======================================================================

\begin{center}
\renewcommand{\arraystretch}{1.25}
\begin{ShadedTabular}{p{5.2cm}p{8.0cm}}
  \toprule
  \textbf{Symptom} & \textbf{Cause} \\
  \emph{``Glossary source \dots{} is in the old (7.10.x) format''} &
    A \texttt{.tex} entry file is still registered.  Run \verb!make glsbib!. \\
  Entries in a strange order &
    Missing \texttt{sort} keys (step~2). \\
  An entry is missing from a list &
    For acronyms this is normal --- only the ones you cite are listed.  For
    the glossary and symbols lists it is not; check the entry really is in
    the \texttt{.bib} file. \\
  \emph{``Undefined control sequence \textbackslash bibglsnew\dots{}''} &
    \texttt{bib2gls} did not run.  Check \verb!bib2gls --version!, i.e.\
    check Java. \\
  \bottomrule
\end{ShadedTabular}
\end{center}

Nothing outside the glossaries changed: your chapters, bibliography, cover
configuration and \latexinline{\ntsetup} options all work as before.
